\section{Discussion}

\subsection{Interpreting the Hybrid Model's Success: Mechanisms and Limitations}

The 14.48\% RMSE improvement achieved by the hybrid GARCH-LSTM model over the baseline reflects two distinct mechanisms. First, the nonlinear residual-learning architecture itself captures state-dependent patterns in volatility not explained by returns alone—the ablation test shows a 14.26\% improvement with no sentiment features. Second, sentiment contributes marginal additional information (0.22\% incremental improvement), suggesting that while news carries some signal, its practical value is constrained by the positive bias pervasive in Vietnamese financial media.

\textit{Why does GARCH-X fail?} The linear specification's failure ($t_{DM} = -6.339$) reveals an important econometric lesson: directly embedding sentiment into the variance equation causes the model to fit spurious training-sample correlations driven by the unconditional positive bias, which breaks out-of-sample when that correlation destabilizes. The LSTM sidesteps this by learning \textit{nonlinear interactions} between sentiment and recent volatility dynamics—for example, "negative news on a calm day signals heightened volatility ahead," whereas "negative news on a volatile day may be merely noise." Such state-dependent effects are invisible to linear models.

\textit{Regime dependence and practical deployment:} The 51\% RMSE improvement in calm regimes versus 15\% degradation in crisis regimes is economically meaningful. The model is most accurate when needed (most trading days are calm) but least accurate when risk is highest (crises). For portfolio managers or risk officers, this suggests a \textit{state-switching strategy}: use the hybrid model for normal-volatility days (improving efficiency and reducing capital overallocation), but switch to or revert toward the pure GARCH baseline during detected crises or high-uncertainty periods to preserve robustness.

\subsection{Connection to Prior Findings: Asymmetry and Sentiment in Emerging Markets}

Our results align with and extend prior work on Vietnamese market dynamics. \citet{vu2023sentiments} documented that negative news significantly increases return variance while positive news does not—a sentiment asymmetry reflecting retail investor behavior. Our nonlinear model implicitly captures this asymmetry: the LSTM learns that negative sentiment shares (low-frequency, high-impact signals) drive volatility more reliably than positive sentiment shares (high-frequency, low-noise background). The feature ablation test confirms this: even without explicit asymmetry features, the LSTM achieves most of the forecasting gain by learning such patterns from the 15-day lagged sequences.

Relative to \citet{bodilsen2025exploiting}, who showed that macroeconomic news sentiment improves HAR-model forecasts on the S\&P 500, we find similar magnitude improvements (16.05\% for Garman-Klass vs. their reported single-digit percentage gains). The larger effect in our Vietnamese context likely reflects higher retail investor share and behavioral volatility in emerging markets, consistent with emerging market microstructure literature.

\subsection{Limitations and Caveats}

Several limitations circumscribe our findings:

\begin{enumerate}
    \item \textbf{Sentiment measurement error:} PhoBERT-based classification introduces noise relative to hand-labeled gold standards. While Sontung (2021) reported 93\% title-only accuracy on CafeF data, our title+lead-paragraph approach may differ. Sentiment scores are proxies, not ground truth, limiting causal interpretation.
    
    \item \textbf{Training data representation:} The LSTM trained on 2015--2021 data did not experience the 2022 liquidity shock, causing poor performance during high-volatility regimes. Periodic retraining or rolling-window approaches may mitigate this, but out-of-distribution tail risk remains a challenge.
    
    \item \textbf{Positive media bias:} Vietnamese financial news exhibits 84\% positive article share, reducing sentiment variation and limiting information content. Macro economic shocks that reverse the positive tone (rare events like the 2022 shock) generate stronger market impacts, but such reversals are undersampled in training.
    
    \item \textbf{No transaction costs:} Our evaluation assumes frictionless trading; real implementation requires considering market-impact costs, trading delays, and liquidity. Forecast improvements sufficient in-sample may not justify trading costs.
    
    \item \textbf{Temporal alignment:} Our trading-day alignment protocol is ad hoc. Articles published around market open (before 14:45) influence day $t$ volatility predictions, but market microstructure effects (bid-ask bounce, opening auctions) may create temporal misalignment not fully captured by our rule.
\end{enumerate}

\subsection{Implications for Market Participants}

\textit{For risk managers:} The hybrid model's superior calm-regime forecasts reduce unnecessary conservatism in position sizing on normal days, improving capital efficiency. However, the regime-dependent performance suggests combining the hybrid forecast with crisis-detection signals (e.g., recent volatility spikes, VIX-style indicators) to revert to more robust baselines during stress.

\textit{For traders:} Volatility forecasts enable more efficient option pricing, straddle selling, and variance-swap trading. The 48\% MAPE improvement translates to more accurate pricing of variance products, potentially exploitable via algorithmic trading if execution costs are sufficiently low.

\textit{For policymakers:} The finding that sentiment predicts volatility in a retail-dominated market confirms that investor behavior (not just fundamental information) drives market dynamics in emerging economies. This supports policies promoting financial literacy and reducing herd-driven volatility, such as enforcing cooling-off periods for retail options trading or implementing position limits.

\section{Conclusion}

This paper develops and evaluates a two-stage hybrid GARCH-LSTM model for forecasting volatility of the VN-Index using Vietnamese financial news sentiment extracted via PhoBERT. Our primary findings are:

\begin{enumerate}
    \item \textbf{Hybrid architecture outperforms linear baselines:} The hybrid model achieves 14.48\% RMSE improvement (79.24\% MAPE reduction) relative to GARCH(1,1), with highly significant Diebold-Mariano tests ($p < 0.001$). Critically, linear GARCH-X specifications fail ($t_{DM} = -6.339$), confirming that nonlinear learning is essential.
    
    \item \textbf{Nonlinearity, not sentiment alone, drives gains:} The ablation test shows 14.26\% improvement without sentiment features, indicating the residual-learning architecture is the primary driver. Sentiment contributes marginal additional value, limited by media positive bias.
    
    \item \textbf{Substantial regime dependence:} The model excels in calm periods (51\% RMSE reduction) but underperforms during crises (15\% degradation), suggesting practical deployment requires state-switching strategies combining the hybrid model with robustness mechanisms.
    
    \item \textbf{Robustness across specifications:} Results are stable under alternative sentiment thresholds, volatility targets, and estimation windows, confirming findings generalize beyond the baseline specification.
\end{enumerate}

\subsection*{Contributions}

This work contributes to financial econometrics and NLP finance research in three ways:

\begin{enumerate}
    \item \textbf{Novel Vietnamese financial dataset:} We construct the first large-scale linked dataset of Vietnamese stock prices and CafeF news with PhoBERT sentiment extraction, enabling future research on emerging-market news effects.
    
    \item \textbf{Empirical validation of hybrid architectures:} We demonstrate on Vietnamese data that hybrid econometric-deep learning models can substantially improve forecasts, supporting the architectural insights of \citet{leber2025nonlinear} in a new market context.
    
    \item \textbf{Practical insights for emerging markets:} We establish that sentiment's volatility impact operates through nonlinear, regime-dependent channels in retail-dominated markets, informing both academic understanding and practitioner strategies.
\end{enumerate}

\subsection*{Future Research}

Future work could address several directions:

\begin{itemize}
    \item \textbf{Hand-labeled validation:} Benchmark PhoBERT sentiment against expert annotations on a CafeF subsample to quantify measurement error.
    \item \textbf{Online learning and adaptation:} Implement online LSTM retraining or ensemble methods to improve crisis-regime performance.
    \item \textbf{Causal mechanisms:} Use structural VARs or nonlinear Granger causality to isolate sentiment's causal contribution versus mere correlation.
    \item \textbf{Multivariate forecasting:} Extend to joint prediction of returns and volatility, testing whether sentiment predicts returns through volatility or independently.
    \item \textbf{Alternative data sources:} Incorporate social media sentiment (Twitter, stock forums) and quantify their incremental value relative to news alone.
\end{itemize}
